\section{Related Work}
\textbf{Skill attacks and benchmark scope.} Skill-Inject treats Skill text and optional scripts as injection surfaces, and distinguishes contextual actions under warning and legitimizing policies~\cite{skillinject2026}. Its native agent scaffolds and tasks are broader than our structured executor. AgentDojo evaluates utility and attacks in stateful tool environments~\cite{agentdojo2024}, while ASB organizes attacks across several agent components~\cite{asb2024}. Their results motivate joint security--utility measurement; their attack rates are not directly comparable to rates from our fixtures. SCR-Bench shows why effects across composed Skill paths require state and path observations~\cite{scrbench2026}. We study one bundle per run and do not reproduce its multi-Skill composition setting.

\textbf{Task-aware permissions.} SkillScope builds a unified graph, validates action over-privilege under instantiated tasks, and projects control-flow constraints back into Skill material~\cite{skillscope2026}. Its runtime validation and rewriting exceed the conservative action-check approximation used here. SkillGuard treats Skills as permission-bearing artifacts and governs both context influence and action effects~\cite{skillguard2026}. Progent represents privileges as symbolic tool-name and argument policies, with checked policy updates~\cite{progent2025}. IntentCap assigns different authority fields to different context sources and constrains capability composition~\cite{intentcap2026}. Our separation of task authority from untrusted Skill requests follows this line of work; it is not a claim to originate contextual least privilege.

\textbf{Semantic graphs and runtime evidence.} Semia reconstructs a relational Skill representation and queries structural risk, while explicitly separating a documented gate from proof that it executes~\cite{semia2026}. SEFZ targets violations of Skill-declared guardrails under benign inputs~\cite{sefz2026}. VIGIL checks finite execution traces against behavioral policies, including event relationships and value flow~\cite{vigil2026}. CaMeL separates trusted control from untrusted data and mediates capabilities~\cite{camel2025}. These methods motivate our distinction between potential graph paths and committed runtime effects. Our bounded backward search and closed-tool provenance are much narrower than general semantic verification or whole-program information-flow analysis.

\textbf{Repair and the comparison gap.} Aletheia tests whether permissions can be withheld while independently specified functionality still passes, interpreting successful restrictions as task-relative dispensability witnesses~\cite{aletheia2026}. That framing cautions against equating a small syntactic edit with a globally minimal policy. Here, repair is minimum-cost only within an enumerated set of sink restrictions; utility is measured after execution. The central comparison is between goal-selected repair and the same filtering power driven solely by a task contract. We therefore report B2R as a strong graph-free comparator and label the SkillScope-inspired condition as an approximation, not an official reproduction or a state-of-the-art ranking.
